\section*{Расчётные данные и воспроизведение}
\label{sec:archive}
Данные организованы относительно каталога проекта
\nolinkurl{phase_reproduction/}. Исходный файл потенциала —
\nolinkurl{inputs/fitted.mtp}, тип 0 соответствует Ti, тип 1 — Al.
Его SHA-256 по реестру запусков и независимой проверке энергии равен
\begin{center}
{\small\ttfamily 704f7fd900c7446286066b741084c30ebc7fcf26e5b9b8fd2f29a18ba4ffd661}.
\end{center}
\begin{itemize}
\item Кластерные функции и энергии. \nolinkurl{atomistic/atat/output/design.csv}
описывает выборку для независимой подгонки. В \nolinkurl{author_comparison/}
файл \nolinkurl{cells/orbit_mapping.csv} сопоставляет орбиты, а
\nolinkurl{eci/matched_ECI.csv} и \nolinkurl{eci/cross_validation.csv}
содержат коэффициенты и результаты переноса между выборками.
\item Решёточное равновесие. В том же каталоге
\nolinkurl{thermo/mfa_1000K.csv} и \nolinkurl{mc/mc_same_windows.csv}
содержат результаты среднего поля и две обработки окон MC.
Рисунок~\ref{fig:ce-temp} использует столбцы
\texttt{author\_alpha} и \texttt{author\_alpha2} для обеих подгонок,
одну процедуру определения границ и наибольшую доступную ячейку
при каждой из показанных температур. Перестановки вкладов и зависимость
от пути сравнения приведены в \nolinkurl{qha/qha_attribution.csv}
и \nolinkurl{qha/qha_swaps.csv}.
\item Малые ячейки. \nolinkurl{alloy_ti/targeted_comparison.json} содержит
оценки с дополнительной выборкой для четырёх наиболее весомых классов;
\nolinkurl{alloy_ti/cage_sensitivity.json} — независимые оценки при двух
ограничениях. Геометрию короткой связи и независимый пересчёт энергии
сохраняет \nolinkurl{cluster/runs/alloy-ti-geometry-01/outputs/}.
Кривые на рис.~\ref{fig:ti-integrand} получены из трёх записей
\texttt{samples\_*.npz} в архиве
\nolinkurl{cluster/runs/alloy-ti-n8-01/outputs/draws.tar.gz}.
Для каждой записи применены её логарифмические веса и условие \(W<0{,}40\).
\item Атомистическая бинодаль. В \nolinkurl{solvus_1000/series/} каждой
серии соответствует каталог с \texttt{windows.csv}, записями измерений и
конечными состояниями. В \nolinkurl{solvus_1000/run_provenance.json}
проверяются номера заданий UMT, хеши полученных архивов и состав серий.
В \nolinkurl{solvus_1000/summary_statistics.json} приведены результаты
для десяти ансамблей и попарные изменения границ.
\item Сравнение с C46. \nolinkurl{solvus_1000/c46_reference.json}
содержит тринадцать исходных пар «состав — химический потенциал»,
результат общей касательной, условные интервалы и хеши авторских файлов.
\end{itemize}

Для повторения основного анализа нужны Python, NumPy и SciPy. Следующие
команды PowerShell запускаются из \nolinkurl{phase_reproduction/solvus_1000/}.
Они восстанавливают обе изотермы для 1400 атомов со свободным \(c/a\) и
их условные статистические интервалы из завершённых траекторий.
{\small
\begin{verbatim}
$m = 's1400w35a','s1400w35b','s1400w40a','s1400w40b' |
     ForEach-Object { "series/$_/windows.csv" }
python analysis.py @m --bootstrap 192 --seed 26713 --output replay.json
\end{verbatim}
}
